% META: {"id":"1SPE-SUITES-EX-031","chapitre":"1SPE-SUITES","type_objet":"exercice","capacites":["1SPE-SUITES-C3","1SPE-SUITES-C5"],"capacites_codes":["C3","C5"],"methodes":["M3","M5"],"parcours":2,"competences":["calculer","raisonner"],"duree_min":20,"mode_creation":"ex_nihilo","sources_inspiration":[],"parametres_sympy":{"u0":{"domaine":"entier","min":1,"max":5},"q":{"domaine":"rationnel","contrainte":"q>1"}},"fichier_tex":"chapitres/1SPE-SUITES/exercices/1SPE-SUITES-EX-031.tex","corrige_tex":"chapitres/1SPE-SUITES/corriges/1SPE-SUITES-CO-031.tex","status":"approved","origin":"generated","status_history":["generated","audited","accepted","approved"],"release_acceptance":"RELEASE_OWNER_BATCH_ACCEPTANCE","acceptance_closure_digest":"sha256:9b3ccf9a81c5520fbb7e03b7b2d3e7bf2057a02834b6d908bf3be5f49f3b553f","acceptance_receipt_digest":"sha256:4fad86f0282e0fa4e0419cca1ad6379ae7af5a280c04a4a90880f90a1c044242"}
% BEGIN-VERIFY
% from sympy import Rational, simplify, symbols
% n = symbols('n', integer=True, nonnegative=True)
% u0 = 2
% q = Rational(5, 3)
% u = u0 * q**n
% assert u.subs(n, 0) == 2
% assert u.subs(n, 1) == Rational(10, 3)
% ratio = simplify(u.subs(n, n+1) / u)
% assert ratio == Rational(5, 3)
% # q > 1 => suite croissante
% assert ratio > 1
% # u_n croissante : u_{n+1}/u_n = 5/3 > 1
% v = 1 / u
% assert simplify(v - Rational(1, 2) * Rational(3, 5)**n) == 0
% assert simplify(v.subs(n, n+1) / v) == Rational(3, 5)
% END-VERIFY
\begin{exercice}{1SPE-SUITES-EX-031}{2}{20}
Soit $(u_n)$ la suite définie pour tout entier naturel $n$ par $u_n = 2 \times \left(\dfrac{5}{3}\right)^n$.

\begin{enumerate}
  \item Reconnaître directement, à partir de sa forme explicite, que $(u_n)$ est une suite géométrique. Préciser sa raison et son premier terme.

  \item Étudier le sens de variation de $(u_n)$ :
  \begin{enumerate}
    \item en calculant et étudiant le quotient $\dfrac{u_{n+1}}{u_n}$ ;
    \item en calculant et étudiant la différence $u_{n+1} - u_n$ (en factorisant le résultat).
  \end{enumerate}
  Vérifier que les deux méthodes donnent la même conclusion.

  \item Pour tout entier naturel $n$, on pose $v_n = \dfrac{1}{u_n}$. Exprimer $v_n$ en fonction de $n$, puis montrer que $(v_n)$ est une suite géométrique. Préciser sa raison et son premier terme.

  \item En déduire le sens de variation de $(v_n)$ sans nouveau calcul, puis comparer ce comportement à celui de $(u_n)$.
\end{enumerate}
\end{exercice}
